\chapter{Módulo 3 --- Programación: control de flujo y funciones}

\begin{itemize}
\tightlist
\item
  \textbf{Duración estimada:} 1.5 h teoría + 2.5 h práctica
\item
  \textbf{Objetivos de aprendizaje.} Al finalizar, la persona podrá:

  \begin{enumerate}
  \def\labelenumi{\arabic{enumi}.}
  \tightlist
  \item
    \textbf{Implementar} decisiones con
    \passthrough{\lstinline!if!}/\passthrough{\lstinline!else!},
    \passthrough{\lstinline!ifelse()!} y
    \passthrough{\lstinline!dplyr::case\_when()!}.
  \item
    \textbf{Ejecutar} iteraciones con \passthrough{\lstinline!for!} y
    \passthrough{\lstinline!while!}.
  \item
    \textbf{Escribir} funciones propias con argumentos, valores por
    defecto y validaciones.
  \item
    \textbf{Comparar} un bucle con la alternativa vectorizada o
    funcional (\passthrough{\lstinline!sapply()!},
    \passthrough{\lstinline!purrr::map()!}).
  \end{enumerate}
\end{itemize}

\section{Contenidos}

\begin{enumerate}
\def\labelenumi{\arabic{enumi}.}
\tightlist
\item
  Operadores lógicos: \passthrough{\lstinline!\&!},
  \passthrough{\lstinline!|!}, \passthrough{\lstinline"!"},
  \passthrough{\lstinline!\&\&!}, \passthrough{\lstinline!||!},
  \passthrough{\lstinline!\%in\%!}.
\item
  \passthrough{\lstinline!if!}, \passthrough{\lstinline!else if!},
  \passthrough{\lstinline!else!}; \passthrough{\lstinline!ifelse()!}
  vectorizado.
\item
  Bucles \passthrough{\lstinline!for!} y
  \passthrough{\lstinline!while!}; \passthrough{\lstinline!break!} y
  \passthrough{\lstinline!next!}.
\item
  Anatomía de una función: nombre, argumentos, cuerpo, valor de retorno.
\item
  Valores por defecto y validación con \passthrough{\lstinline!stop()!}.
\item
  Alcance (scope) de variables.
\item
  Vectorización y familia \passthrough{\lstinline!apply!}
  (\passthrough{\lstinline!sapply!}, \passthrough{\lstinline!lapply!})
  vs.~\passthrough{\lstinline!purrr::map()!}.
\end{enumerate}

\section{Desarrollo}

\subsection{Condicionales}

\begin{lstlisting}[language=R]
unidades <- 25
if (unidades >= 20) {
  "Alta rotación"
} else if (unidades >= 10) {
  "Rotación media"
} else {
  "Baja rotación"
}

u <- c(3, 12, 25)
ifelse(u >= 10, "Suficiente", "Reabastecer")   # vectorizado
"Tablet" %in% c("Tablet", "Reloj")             # TRUE
\end{lstlisting}

\subsection{Bucles}

\begin{lstlisting}[language=R]
productos <- c("Smartphone A", "Tablet", "Reloj")
for (p in productos) {
  print(paste("Revisando inventario de", p))
}

stock <- 50
semanas <- 0
while (stock > 0) {
  stock <- stock - 12
  semanas <- semanas + 1
}
semanas        # semanas hasta agotar: 5
\end{lstlisting}

\subsection{Funciones}

\begin{lstlisting}[language=R]
ingreso <- function(precio, unidades, descuento = 0) {
  if (descuento < 0 || descuento >= 1) {
    stop("`descuento` debe estar en [0, 1).")
  }
  precio * unidades * (1 - descuento)
}

ingreso(4999, 10)
ingreso(4999, 10, descuento = 0.15)
ingreso(c(899, 2499), c(20, 8))     # funciona con vectores
\end{lstlisting}

\subsection{Vectorizar en lugar de iterar}

\begin{lstlisting}[language=R]
x <- 1:5
# Con bucle
cuadrados <- numeric(length(x))
for (i in seq_along(x)) cuadrados[i] <- x[i]^2
cuadrados
# Vectorizado (preferido)
x^2
# Funcional
sapply(x, function(v) v^2)
library(purrr)
map_dbl(x, \(v) v^2)               # \(v) es la sintaxis corta de function(v), R >= 4.1
\end{lstlisting}

\section{Actividades y ejercicios}

\begin{enumerate}
\def\labelenumi{\arabic{enumi}.}
\tightlist
\item
  Escribe \passthrough{\lstinline!clasificar\_rotacion(u)!} que devuelva
  \passthrough{\lstinline!"Alta"!}, \passthrough{\lstinline!"Media"!} o
  \passthrough{\lstinline!"Baja"!} para un \textbf{vector} de unidades
  (umbrales 20 y 10).
\item
  Escribe \passthrough{\lstinline!margen(precio, costo)!} que devuelva
  el margen porcentual y lance error si
  \passthrough{\lstinline!costo > precio!}.
\item
  Con un bucle \passthrough{\lstinline!for!}, imprime el total de
  unidades por producto en \passthrough{\lstinline!ventas!}. Luego obtén
  lo mismo con \passthrough{\lstinline!tapply()!}.
\item
  Simula lanzar un dado 1 000 veces con
  \passthrough{\lstinline!sample()!} y calcula la proporción de seises.
\end{enumerate}

\subsection{Soluciones}

\begin{lstlisting}[language=R]
# 1 (dplyr::case_when es vectorizado y legible)
clasificar_rotacion <- function(u) {
  dplyr::case_when(
    u >= 20 ~ "Alta",
    u >= 10 ~ "Media",
    TRUE    ~ "Baja"
  )
}
clasificar_rotacion(c(3, 12, 25))

# 2
margen <- function(precio, costo) {
  if (any(costo > precio)) stop("El costo no puede ser mayor que el precio.")
  (precio - costo) / precio * 100
}
margen(100, 60)                   # 40

# 3
ventas <- read.csv("datos/ventas.csv")
for (p in unique(ventas$producto)) {
  total <- sum(ventas$unidades[ventas$producto == p], na.rm = TRUE)
  cat(p, ":", total, "\n")
}
tapply(ventas$unidades, ventas$producto, sum, na.rm = TRUE)

# 4
set.seed(1)
dado <- sample(1:6, 1000, replace = TRUE)
mean(dado == 6)
\end{lstlisting}

\section{Evaluación (sumativa parcial)}

Entrega un script \passthrough{\lstinline!funciones\_ventas.R!} con 3
funciones documentadas (comentario de propósito, argumentos y retorno) y
al menos una validación con \passthrough{\lstinline!stop()!}. Rúbrica:

\begin{longtable}[]{@{}
  >{\raggedright\arraybackslash}p{(\columnwidth - 6\tabcolsep) * \real{0.2500}}
  >{\raggedright\arraybackslash}p{(\columnwidth - 6\tabcolsep) * \real{0.2500}}
  >{\raggedright\arraybackslash}p{(\columnwidth - 6\tabcolsep) * \real{0.2500}}
  >{\raggedright\arraybackslash}p{(\columnwidth - 6\tabcolsep) * \real{0.2500}}@{}}
\toprule\noalign{}
\begin{minipage}[b]{\linewidth}\raggedright
Criterio
\end{minipage} & \begin{minipage}[b]{\linewidth}\raggedright
3 --- Logrado
\end{minipage} & \begin{minipage}[b]{\linewidth}\raggedright
2 --- En proceso
\end{minipage} & \begin{minipage}[b]{\linewidth}\raggedright
1 --- Inicial
\end{minipage} \\
\midrule\noalign{}
\endhead
\bottomrule\noalign{}
\endlastfoot
Correctitud & Las 3 funciones dan el resultado esperado & 2 funciones
correctas & ≤ 1 correcta \\
Vectorización & Aceptan vectores sin bucles innecesarios & Parcial &
Solo escalares \\
Validación & Mensajes de error claros & Validación incompleta & Sin
validación \\
Estilo & Nombres claros, comentarios útiles & Aceptable & Difícil de
leer \\
\end{longtable}

\section{Referencias verificadas}

\begin{itemize}
\tightlist
\item
  \textbf{Libro:} Wickham, H., Çetinkaya-Rundel, M., \& Grolemund, G.
  (2023). \emph{R for Data Science} (2.ª ed.). O'Reilly Media. ISBN
  978-1-4920-9740-2. Español: \url{https://davidrsch.github.io/r4ds-es/}
  {[}ES{]} {[}OA{]} --- caps. ``Funciones'' e ``Iteración''.
\item
  \textbf{Libro:} Grolemund, G. (2014). \emph{Hands-On Programming with
  R}. O'Reilly Media. ISBN 978-1-4493-5901-0.
  \url{https://rstudio-education.github.io/hopr/} {[}EN{]} {[}OA en
  línea{]} --- proyecto ``Slot machine'' (if/else, bucles, funciones).
\item
  \textbf{Libro:} Wickham, H. (2019). \emph{Advanced R} (2.ª ed.). CRC
  Press. ISBN 978-0-8153-8457-1. \url{https://adv-r.hadley.nz/} {[}EN{]}
  {[}OA en línea{]} --- cap. 5 ``Control flow'' y cap. 6 ``Functions''.
\end{itemize}
